\cotefile{10}
\scanpage{123}
% archive sheet number: 122
% source lines 1--20
$\handwritten{\pi}_i^{\handwritten{2}}=\handwritten{\pi}_i,\handwritten{\pi}_i\handwritten{\pi}_j=0$ pour $i\handwritten{\ne}j$, $\handwritten{\sum_i\pi}_i=1$. La donnée d'une $\otimes$-structure sur $C$ impliquant les données a)b),et celles-ci compatibles avec le foncteur $C\handwritten{\to}D$ (pour $c=\mathrm{id}$),équivaut alors à une structure d'algèbre associative unitaire sur $A$, pompatible avec la structure de coalgèbre, et pour laquelle $A\handwritten{\to}k\handwritten{[}T,T^{-1}\handwritten{]}$ est un homomorphisme d'\unsure{algèbres|Àlgèbres}. Comment s'exprime le fait que \deleted{\unsure{$C\to D$ est}} sur $C$ il y a une structure de symmétrie \deleted{telle} respectée par le foncteur $C\handwritten{\to}D$ (et $c$ choisi plus haut) ? On peut d'abord exprimer la structure\inserted{\handwritten{$\otimes$}} précédente par une application diagonale de $U$ satisfaisant $\handwritten{\Delta}(\handwritten{\pi}_i)=\handwritten{\sum}_{j+k=i}\handwritten{\pi}_j\otimes\handwritten{\pi}_k$. \unsure{La|Ea} condition cherchée \deleted{revient à dire que si $M$ est un $U$-module de degré $i$, \unsure{N.e.|i.e.} tel que ${}_iM=M$, et $N$ de degré $j$ i.e. ${}_jN=N$, alors \unsure{làisomorphisme|laisomorphisme} $x\otimes y\quad(-1)^{ij}y\otimes x$ de $M\otimes_kN$ sur $N\otimes_kM$ est compatible avec les opérations de $U$} est celle-ci: si nous graduons $U$ par ``degré des opérateurs dans des gradués'' i.e. \deleted{$U^i\unsure{=\pi_j|=\sum_j}$} ces quotients $U_{\handwritten{\alpha}}$ par $U_{\handwritten{\alpha}}^i=\handwritten{\widehat{\sum}\pi}_{i+j}U_{\handwritten{\alpha}}\handwritten{\pi}_i$, alors l'application diagonale devient \uline{anti}commutative i.e. la \unsure{muttiplication|multiplication} de $A$ \deleted{\unsure{respecte}} est \uline{anticommutative}. (Il y aurait lieu de prouver que la graduation de $A$\inserted{i.e. de $U$} permet de récupérer le \deleted{\unsure{foncteur au|foncteur}} graduation sur le foncteur oubli,\margin{à vérifier, \handwritten{(\unsure{groupe} \unsure{de|des} classes \unsure{quelc.|quelq.} pour \unsure{graduation} ?)}} i.e. les $\handwritten{\pi}_i\handwritten{\in}U$ i.e. $A\handwritten{\to}k\handwritten{[}T,T^{-1}\handwritten{]}$, et que la graduation de $A$ \deleted{n'est} peut être choisie \unsure{arbitrairement}, avec la seule condition, de toutes façons nécessaire\deleted{s}, que $A$ soit anticommutative et l'homomorphisme diagonal \handwritten{\unsure{soit}} compatible aux graduations).

% source lines 21--27
\handwritten{9/} Quelles que soient les structures a)\deleted{\unsure{b)}}c) en jeu\inserted{(b) -l'unité - étant \handwritten{\unsure{fixée}}}, le groupe des $\otimes$-automorphismes du foncteur \deleted{fibre} ``oubli'' en est indépendant, c'est le groupe des $g\handwritten{\in}U$ tels que $\handwritten{\Delta}(g)=g\otimes g$ et d'augmentation 1\inserted{\handwritten{(i.e. les hom unitaires $\unsure{g|q}:A\to k$)}}\unsure{(}NB cette dernière condition exprime la condition b) de 3º). Remplaçant $k$ par un e algèbre quelconque, on trouve encore un \deleted{\unsure{foncte|foncteur}} schéma en monoïdes affine sur $k$, mais qui en général ne permettra pas de réconstitu\deleted{\unsure{t}}er la situation entièrement. \deleted{\unsure{(donner exemple}} Il faut analyser de plus près
\scanpage{124}
% archive sheet number: 123
% source lines 1--18
le foncteur compatible aux $\otimes$ associatifs unitaires qu'on obtient de $C$ dans $\mathrm{Rep}_k(G)$, et voir en particulier sous quelles conditions il est une équivalence de catégories: \deleted{Alors}\inserted{\handwritten{il \unsure{f. et s.|f. et p.} que $A$ soit commutatif! ``\unsure{Géométriquement|biométriquement}''}} Lorsqu'il en est ainsi,\inserted{\handwritten{\uuline{dans le cas gradué anticomm.}}} la situation est décrite entièrement par la donnée de $G$, de la graduation sur \deleted{\unsure{son}} le foncteur oubli sur \deleted{$C$}\inserted{$R_k(G)$}, décrite par un homomorphisme de monoïdes $\underline{G}_m\handwritten{\to}G$, et enfin la condition que pour deux représentation $M,N$ de $G$, l'isomorphisme de Koszul $M\otimes N\handwritten{\xrightarrow{\sim}}N\otimes M$ sur les vectoriels sous-jacents est compatible avec l'action de $G$. Bien sûr, si $k$ est de car 2, ce n'est pas une restriction, et dans ce cas d'ailleurs Koko ne change rien i.e. tous nos ennuis et distinguos disparaissent. \deleted{Si} Il suffit en tous cas que l'homomorphisme induit $\handwritten{\mu}_2\handwritten{\longrightarrow}G$ soit central;\inserted{\handwritten{et on voit que cela est \unsure{vérifié|vrai} si \unsure{(car $k\ne2$, i.e.)|(car $k\ne2$, c.à.d.)} \unsure{que cela implique} $A^{\unsure{*|\#}}=A^{\unsure{\oplus|0}}$ donc $A^*$ commutatif !}} à fortiori\inserted{\handwritten{il suffit}} que $\underline{G}_m\handwritten{\to}G$ soit central\handwritten{.}\deleted{,i.e.}\inserted{Cela signifie} que la décomposition d'un $M$ en composants \deleted{homogènes}\inserted{pairs et impairs} soit invariante sous $G$ i.e. se fasse dans $C$,resp. qu'il en soit de même pour la décomposition en degrés\inserted{\handwritten{entiers}}. Noter que cette dernière condition s'exprime directement en termes de $C$ (muni de sa structure $\otimes$ et du foncteur \deleted{oubli} gradué\deleted{\unsure{x|s}} $C\handwritten{\to}\mathrm{Grdf}(k)$), \deleted{et on \deleted{espère}\inserted{\handwritten{voit}} que cette condition est suffisante (peut-être nécessaire et suffisante)\inserted{\handwritten{si car. $k\ne2$}} pour qu'on soit dans le cas favorable\inserted{\handwritten{où}} $C\handwritten{\to}\mathrm{Rep}_k(G)$ \unsure{est} une équivalence.}
\margin{\deleted{\handwritten{et on \unsure{voit}}}}

% source lines 19--27
\handwritten{10)} Ce résultat s'appliquera alors, en particulier, à la catégorie des motifs\inserted{semi-simples} (car les \deleted{foncteurs} graduations naturelles sur les foncteurs $H^\bullet$ qu'on arrive à y définir \deleted{sont} proviennent d'une graduation de la catégorie des motifs, - ce qui n'est qu'une traduction du fait que les correspondances\inserted{cohomologiques} de Kunneth sur une variété algébrique projective lisse sont algébriques). On respire ! Mais on fera attention que les foncteurs $H^\bullet$ en question ne sont pas des foncteurs fibres proprement, mais des foncteurs fibres tordus.(On ne pourra s'empêcher de les appeler quand-même foncteurs fibres!) Quand on travaille avec des motifs pas nécessairement
\scanpage{125}
% archive sheet number: 124
% source lines 1--11
semi-simples, on a une catégorie qui n'est plus graduée, seulement filtrée, \deleted{\unsure{I}} Si $V$ est la catégorie des schémas de type fini sur un corps fixé, $M$ la catégorie des motifs sur ce même corps (\unsure{présumée|pvrésumée} avec toutes les données et conditions de 1º)\handwritten{\deleted{\unsure{(}}}, on doit avoir un foncteur \deleted{\unsure{$V$|$Y$}} compatible avec ces structures \deleted{de} (cohomologie motivique)
\[
h:V\handwritten{\longrightarrow}\unsure{\mathbb{G}|G|E}\mathrm{rad}(M)\handwritten{;}\quad\text{(structure de Koko sur }\mathrm{Grad}(M)\text{)},
\]
et on désire que les foncteurs cohomologiques habituels sur $V$ se factorisant par \deleted{$M$. \unsure{On impose}}\inserted{des foncteurs sur $M$. On voudra} que les foncteurs \deleted{fibres} obtenus sur $M$ \unsure{soinnt|sointnt|soient} des foncteurs fibres \deleted{\uline{tordus} seulement} véritables\inserted{\handwritten{(non tordus)}}, ce qui n'est pas idiot, \deleted{car}\inserted{\handwritten{car}} on avait mis la structure de symétrie de Koko sur $\mathrm{Grad}(M)$: on s'est débarassé ainsi une fois pour toutes de la ``torsion''.

% source lines 12--15
On peut aussi, pour parler de vrais foncteurs fibres sur la catégorie des motifs semi-simples, remplacer la donnée de symmétrie ``géométrique'' i.e. de Koko pour le produit tensoriel dans \deleted{$\mathrm{Grad}(M)$}\inserted{la catégorie graduée $M$}, par la même canulée par le signe $(-1)^{ij}$ quand on a des \unsure{arguments} de degré $i,j$.

% source lines 16--24
11) Bien sûr, il y a lieu de chercher des théorèmes de structure intrinsèques, en termes de gerbes, pour des catégories $\handwritten{\unsure{\mathscr{C}|\mathfrak{C}|\mathbb{C}}}$ munies des structures de 1º,\inserted{avec $k=\mathrm{End}(\underline{1})$ un corps donné.} Il y a lieu au besoin de supposer qu'il existe une extension de $k$ sur laquelle il y ait un foncteur fibre \unsure{tordu|toodu} gradué (ou mltigradué ? Mais où s'arrêter ? ). Le plus beau serait évidemment de trouver un théorème de structure \unsure{ggnéral|ggénéral|général}, sans \uline{supposer} l'existence d'un foncteur fibre plus ou moins tordu où que ce soit, en définissant une notion assez générale de foncteur fibre tordu (?) et en ppuvant qu'\unsure{ip|il} y en a toujours sur une extension convenable de $k$.
\scanpage{126}
% archive sheet number: 125
% source lines 1--9
Voici quelques \unsure{cogitations|cohitations} dans ce sens, lorsque $C$ est muni d'un foncteur additif exact fidèle\deleted{\unsure{F}} $k$-linéaire $F:C\handwritten{\to}\unsure{\mathbb{M}|M}\mathrm{odf}(k)$, et \unsure{correspond} à ce titre à une algèbre profinie (ass un.) $U$, resp. une coalgèbre ass. unit. $A$, sur le corps $k$ (cf nº6). On \deleted{voit que la donnée d} suppose donné une opération $\otimes$, compatible avec $F$, i.e. un\deleted{\unsure{e}} hom. diagonal $U\handwritten{\to}U\overset{\handwritten{\scriptscriptstyle\wedge}}{\otimes}U$ i.e. une composition $A\otimes_kA\handwritten{\to}A$. On se propose de déterminer\inserted{\uline{toutes}} les données de symétrie $\handwritten{\psi}$ pour $\otimes$ (\uline{pas seulement} celles compatibles avec le foncteur $F$ et la donnée de symmétrie ordinaire sur $\mathrm{Modf}(k)$ !) On trouve qu'\unsure{illes|elles} correspondent aux
\margin{\handwritten{[Dans le cas du nº 6, \uline{\unsure{gr. deux|gr. deu.}}, on aura $\sigma=\pi^+\otimes\pi^++\pi^+\otimes\pi^-+\pi^-\otimes\pi^+-\pi^-\otimes\pi^-$]}}
\[
\handwritten{\sigma\in}(U\overset{\handwritten{\scriptscriptstyle\wedge}}{\otimes}U)^{\handwritten{*}}
\tag*{\handwritten{(0)}}
\]
satisfaisant les conditions
\[
\left\{\begin{aligned}
&\handwritten{\sigma'\sigma}=1\quad\text{(i.e. }\handwritten{\sigma}\text{ et }\handwritten{\sigma'}\text{ inverses l'un de l'autre: on aura aussi}\\
&\qquad\handwritten{\sigma\sigma'}=1,\quad\text{puisque }\handwritten{\sigma}\text{ est supposé inversible)}\\
&\boxed{\handwritten{\sigma(\Delta}u)\handwritten{\sigma}^{-1}=\handwritten{(\Delta}u\handwritten{)'}\quad\text{pour tout }u\handwritten{\in}U}
\end{aligned}\right.
\tag*{\handwritten{(1)}}
\]
\margin{\handwritten{condition de symétrie ``tordue'' de l'application diagonale.}}
% source lines 14--18
en désignant par $x\handwritten{\mapsto}x'$ l'automorphisme de symétrie de $U\overset{\handwritten{\scriptscriptstyle\wedge}}{\otimes}U$. NB On ne dit pas que $\handwritten{\Delta}U$ est stable sous $x\handwritten{\leftrightarrow}x'$, mais on considère $\handwritten{\Delta}$ et le symétrique $\handwritten{\Delta'}$ défini par $\handwritten{\Delta'}(u)=\handwritten{\Delta}(u)\handwritten{'}$, qui définit la loi $\otimes$ \uline{opposée} de celle définie par $\handwritten{\Delta}$, et les relations entre $\handwritten{\sigma}$ et $\handwritten{\sigma'}$ sont symétriques et s'expriment\inserted{\handwritten{en disant}} que $\handwritten{\sigma}$ transforme $\handwritten{\Delta}$ en $\handwritten{\Delta'}$, $\handwritten{\sigma'}$ transforme $\handwritten{\Delta'}$ en $\handwritten{\Delta}$.

% source lines 19--22
Lorsqu'il y a une unité bilatère pour $\otimes$ et que $F$ est compatible avec celle-ci,\inserted{i.e. $A$ admet une unité bilatère,} alors la compatibilité de $\handwritten{\psi}$ à l'unité de $C$ s'exprime par les relations
\[
(\mathrm{id}_U\otimes\handwritten{\varepsilon}_U)(\handwritten{\sigma})=(\handwritten{\varepsilon}_U\otimes\mathrm{id}_U)(\handwritten{\sigma})=1,\quad\text{i.e.}\quad\handwritten{\varepsilon\in}1+U^{\handwritten{+}}\overset{\handwritten{\scriptscriptstyle\wedge}}{\otimes}U^{\handwritten{+}}.
\tag*{\handwritten{(2)}}
\]

% source lines 23--26 and bottom handwritten addition
Lorsqu'on a une donnée d'associativité \deleted{sur} $\handwritten{\phi}$ pour $\unsure{\otimes|\varpi}$, compatible avec $F$ et l'associativité ordinaire de $\otimes$ sur $\mathrm{Modf}(k)$, i.e. $A$ est associative, alors la compatibilité entre $\handwritten{\phi}$ et $\handwritten{\psi}$ s'exprime par une relation du style
\[
\handwritten{\deleted{\unsure{\sigma}_1\unsure{\sigma}_2\unsure{\sigma}_3\unsure{\sigma'}_1\unsure{\sigma'}_2\unsure{\sigma'}_3=1\unsure{\sigma}}\quad(\mathrm{id}_U\widehat{\otimes}\Delta)(\sigma)\deleted{\unsure{\sigma^{23}\sigma^{12}|\sigma^{12}\sigma^{13}}}=\deleted{1}\sigma^{12}\sigma^{13}}
\tag*{\handwritten{(3)}}
\]
\handwritten{où les $\sigma^{ij}$\inserted{\unsure{divers}} sont les 3 images de $\sigma$ dans $U\otimes U\otimes U$ par $x\otimes y\mapsto x\otimes y\otimes1$ et les deux autres appl. analogues.}
\scanpage{127}
% archive sheet number: 126
% source lines 1--4
\deleted{où les $\handwritten{\sigma}_i$ (resp. $\handwritten{\sigma'}_i$) sont déduits de\deleted{s} $\handwritten{\sigma}$ (resp. de $\handwritten{\sigma'}$) par les trois homomorphismes $\handwritten{\delta}_i:U\overset{\handwritten{\scriptscriptstyle\wedge}}{\otimes}_kU\handwritten{\longrightarrow}U\overset{\handwritten{\scriptscriptstyle\wedge}}{\otimes}_kU\overset{\handwritten{\scriptscriptstyle\wedge}}{\otimes}_kU$ \deleted{utilisant} $\handwritten{\Delta}\otimes\mathrm{id}_U$ et symétriques.} (NB je n'ai pas écrit le calcul et ne garantis pas la forme exacte de la relation de l'hexagone).

% source lines 5--13
Soit $k'$ une algèbre (comm. ass. unitaire !) sur $k$, d'où $\unsure{\mathbb{U}|U}'=(U'_{\handwritten{\alpha}})$ et $\unsure{\mathbb{A}|A}'=\underset{\handwritten{\alpha}}{\underset{\handwritten{\longrightarrow}}{\lim}}A'_{\handwritten{\alpha}}$ sur $k'$, avec leurs structures de\deleted{\unsure{algèbres}} bigèbres \handwritten{(}resp. bi-progèbre pour $U'$\handwritten{)}. On se propose de déterminer les foncteurs $G:C\handwritten{\to}\deleted{\mathrm{Modtf}(}\mathrm{Projf}(k')$\inserted{(modules proj de t f sur $k'$)} \uline{exacts et} $k$-\uline{linéaires}, munis d'un isomorphisme $c:G(X\otimes Y)\handwritten{\simeq}G(X)\otimes G(Y)$. \deleted{\unsure{son}} Tout d'abord la donnée de $G^{\unsure{k}}$ exact à droite\inserted{$C\handwritten{\to}\mathrm{Mod}(k')$} revient à la donnée d'un $L=(L_{\handwritten{\alpha}})$, module\inserted{à droite} sur $U$\handwritten{,} muni d'une structure de $k'$-module \deleted{sur $k$} commutant aux opérations de $U$ (NB $L_{\handwritten{\alpha}}=G(U_{\handwritten{\alpha}\unsure{S|s}})$ ), par
\[
G(\mathbb{M})=L\otimes_U\mathbb{M}.
\]

% source lines 14--23
La $k^{\deleted{\unsure{'}}}$-linéarite signifie que les deux structures sur $k$ de $L$ coïncident, i.e. que $L$ provient d'un\deleted{\unsure{e}} module (pro) sur $U'=(U'_{\handwritten{\alpha}})$ (chaque $L_{\handwritten{\alpha}}$ est un module sur $U'_{\handwritten{\alpha}}$).\inserted{On peut alors écrire $G(X)=L\otimes_{U'}X'$, où $X'=X\otimes_kk'$.} Le fait que l'on s'envoie dans $\mathrm{Projf}(k')$ \deleted{s'exprime} et que le foncteur envisagé $G$ est exact s'exprime alors par le fait que les $L_{\handwritten{\alpha}}$ sont \deleted{libres}\inserted{proj.} de type fini \uline{sur} $k'$, et que les $L_{\handwritten{\unsure{\mu|\alpha}}}$ sont \unsure{plets|plats} sur \uline{les} $U_{\handwritten{\alpha}}$ (nécessité triviale, suffisance un petit argument par réduction au cas $X$ simple, puis $U$ un corps égal à $X$ \ldots). Donc on trouve exactement les \uline{modules} $L=(L_{\handwritten{\alpha}})$ \uline{sur} $U'$ \uline{tels que les} $\handwritten{L_\alpha}$ \uline{soient plats sur} $U_{\handwritten{\alpha}}$ \uline{et proj de tf sur} $k'$.
\margin{séparer les conditions, commencer par exactitude = platitude\handwritten{sur les $U_\alpha$}}
D'autre part, indépendamment de ces \deleted{C}\inserted{\handwritten{\unsure{c}}}onditions de platitude,\inserted{\handwritten{si dans $C$ il y a $1$ et $\vee$ \unsure{comp.tibles} avec\inserted{(\unsure{contravariété|contravariants})} $F$, i.e. $\Delta_U$ unitaire et $U$ avec antipodisme}} la donnée d'un isomorphisme $c$ revient à la donnée d'un isomorphisme de bi-$U'$-modules\inserted{\handwritten{à droite (*)}}
\margin{\deleted{\handwritten{\unsure{en le dualisant} \unsure{faisant intervenir} \illegible{2} d'un \unsure{antipodisme} $U'$}}}
\[
L\overset{\handwritten{\scriptscriptstyle\wedge}}{\otimes}_{\handwritten{k'}}L\handwritten{\xleftarrow{\sim}}L\overset{\handwritten{\scriptscriptstyle\wedge}}{\otimes}_{\handwritten{k'}}U'
\quad(\handwritten{\simeq}L\otimes_{(U,\handwritten{\Delta})}(U\overset{\handwritten{\scriptscriptstyle\wedge}}{\otimes}_kU)\deleted{\handwritten{\unsure{\simeq}}}\handwritten{\simeq}L\otimes_{(U',\handwritten{\Delta'})}(U'\overset{\handwritten{\scriptscriptstyle\wedge}}{\otimes}_kU')).
\]
\inserted{\handwritten{\unsure{[|(}\unsure{grâce|voir} à l'antipodisme:\unsure{)|]} $x\otimes u\longrightarrow x\otimes(u\otimes1)$]}}

% source lines 24--26 and lower-margin notes
\deleted{En termes} Utilisant les conditions de platitude, on peut définir $B=\mathrm{Homcont}_{k'}(L,k')$, qui est limite inductive des $B_{\handwritten{\alpha}}$ proj de tf sur $k'$ (à fortiori\inserted{\handwritten{il est}} plat sur $k'$), et une coargèbre sur $A'$ (limite de coalgèbres
\margin{\handwritten{Commençons \unsure{par} tout exprimer du côté de $U,L$, et ne dualisons qu'à la fin}}
\handwritten{(*) La loi de bim. dans $L\otimes_{\unsure{k'|U'}}U'$ est donnée par}
\[
\handwritten{\left\{\begin{aligned}
&\text{1ère loi : à droite des }\unsure{\text{mult.}}\text{ de }U'\\
&\text{2ème loi : donnée par celle de mod : droit}\\
&\text{sur }U\otimes U^o\quad[U^o\text{ opère à droite sur }U'\text{ !}]\\
&\text{via }(\mathrm{id}_U\otimes a.)(\Delta_U)
\end{aligned}\right.}
\]
\scanpage{128}
% archive sheet number: 127
% source lines 1--10
$B_{\handwritten{\alpha}}$ sur les $A_{\handwritten{\alpha}}^{\handwritten{'}}$ qui se réduisent les unes en les autres \ldots), \deleted{\unsure{les $B$}} de plus chaque $B_{\handwritten{\alpha}}$ étant coplat sur $A_{\handwritten{\alpha}}$. La structure\inserted{\handwritten{$c$}} précédente revient à un isomorphisme de \uline{\deleted{bi}\inserted{co}modules} sur \deleted{$\unsure{B'|B|A'}\otimes_{\unsure{k'|k}}$} $A'$
\[
B\otimes_{k'}B\handwritten{\xleftarrow{\sim}}B\otimes_{k'}A'.\quad\deleted{\text{De}}
\tag*{(*)}
\]
Utilisant l'augmentation du facteur $A'$, on en conclut un homomorphisme
\[
B\otimes_{k'}B\handwritten{\longrightarrow}B
\tag*{(**)}
\]
qui sera un \uline{homomorphisme de comodules} sur $A'$ (NB utiliser la multiplication de $A'$ pour définir la structure de $A'$-comodule sur le produit tensoriel sur $k$ de deux tels), i.e. on a une loi de composition sur $B$ ``compatible avec la structure de $A^{\unsure{\mathrm{L}|'}}$-comodule''.
\margin{NB \unsure{Cette|Cettx} équivaut à dire que $B\handwritten{\to}B\otimes_{\unsure{k'|k}}A'$ est hom. d'alg.}
\margin{A préciser}
\margin{\handwritten{$\Vert$}Sauf erreur, la donnée de (*) \uline{équivaut à celle de} (**)}

% source lines 11--27
Si $\handwritten{\Delta}$ sur $U$ est unitaire i.e. $A$ est unitaire i.e. le $\otimes$ sur $X$ \deleted{\unsure{est adme}\illegible{1}} admet une unité bilatère compatible avec $F$, alors la compatibilité de $c$ avec les unités s'exprime, sauf erreur, par la condition que la loi de composition (**) est unitaire.\deleted{\unsure{(}} En effet, il faut écrire que l'on a un isomorphisme \deleted{$\unsure{L}\otimes_{\handwritten{U}}k$} $L\otimes_{U'}k'\handwritten{\simeq}k'$ satisfaisant certaines conditions, or cela revient à se donner une\deleted{élément de} augmentation de $L$ satisfaisant une certaine condition (savoir \deleted{être compatible avec}\inserted{être unité bilatère pour} la structure diagonale de $L$) ou encore un élément de $B$ qui es t une unité bilatère.\inserted{Ceci implique aussi que $B$ est \uline{fidèlement} plat sur $k'$.} De même, si $\handwritten{\Delta}$ est associatif, i.e.\inserted{si} on a pour $\otimes$ une \unsure{dxonnée|donnée} d'associatilité $\handwritten{\phi}$ compatible avec $F$, alors $c$ est compatible avec celle-ci sss $B$ est une algèbre associative. \deleted{Quand les deux conditions sont \unsure{vérifi|vérifié}} Enfin, quand $\handwritten{\Delta}$ est commutative i.e. $C$ avec $\handwritten{\psi}$ donnée de symétrie, \ldots\ alors $c$ est compatible avec cette dernière, et la donnée de symétrie habituelle sur $\unsure{\mathrm{Poojf}|\mathrm{Projf}}(k')$ sss la loi de composition de $B$ est \deleted{\unsure{asso}} commutative. \handwritten{[}(Mais alors, en termes du monoïde affine $\handwritten{\unsure{H|G|\mathscr{H}|\mathscr{G}}}$ défini par $A$ sur $k$, on voit que les structures de $B$ s'expriment en termes de\inserted{\handwritten{$P:$}} $\mathrm{Spec}\ B$ en disant que $\handwritten{H_{k'}}$ y opère à droite,\inserted{\handwritten{et $P$ est \unsure{form.|fibr.} princ. hom.}} et on voit de plus qu'après le changement
\margin{\handwritten{voir si $\unsure{C|G}$ un groupe i.e. si $\exists$ antipodisme.}}
\scanpage{129}
% archive sheet number: 128
% top marginal question
\margin{\handwritten{\uline{NB}} \uline{Question} Sous réserve $B$ associative unitaire, si $A$ est commutative, $B$ l'est-elle aussi ? (On suppose $A$ associative unitaire). Question analogue quand on a des structures graduées avec anticommutativité. Donc, la condition de \deleted{\unsure{symé}} comp. aux sym. pour un $c$ est-elle surabondante, en prés. des autres \unsure{compat.|compa\handwritten{t.}} ?}
% source lines 1--3
de base fidèlement plat $k'\handwritten{\longrightarrow}B$, \deleted{\unsure{le foncteur|le torseur}} il devient trivial, donc on a un ``torseur'' sous le monoïde $\unsure{\mathbb{H}|\mathbb{G}}_{k'}$ ; alors le nouveau foncteur fibre $G$ est déduit de $F'=F\otimes_k k'$ par torsion à l'aide de $G$ \ldots

% source lines 4--14
Mais supposons seu\deleted{\unsure{x}}lement $A$ algèbre associative unitaire, et $B$ itou, sans condition de commutativité.\inserted{On suppose $A$ avec \uline{antipodisme}\inserted{\handwritten{$\underline{a}$}} (s'en passer ?).} Donnons-nous un $\handwritten{\sigma\in}U\otimes_k U$ \deleted{\unsure{ci-dessus}} comme ci-dessus pour introduire une relation de commutativité $\handwritten{\psi}$ . Notons (qu'à cela près que $k'$ n'est plus supposé un corps) on trouve une catégorie $C'=C\otimes_k k'$ sur $k'$, avec $\otimes$, définie\inserted{\handwritten{(soit par $U'$, \uuline{soit} \unsure{dual.|aussi.})}} par le ``proanneau'' $V=G(U)=L\overset{\handwritten{\scriptscriptstyle\wedge}}{\otimes}_U(U,\mathrm{int})$, où $(U,\mathrm{int})$ désigne $U$ muni de la structure de $U$-module \deleted{\unsure{(donc de pro-objet de}} (et même d'objet``pro-anneau'' (pour $\otimes$) de $C$) défini à l'aide de sa structure de bimodule sur $U^o\overset{\handwritten{\scriptscriptstyle\wedge}}{\otimes} U$ \deleted{\unsure{à l'aide de}} par restriction des scalaires via $U\handwritten{\xrightarrow{\Delta}}U\overset{\handwritten{\scriptscriptstyle\wedge}}{\otimes} U\handwritten{\xrightarrow{a\otimes\mathrm{id}_U}}U^o\overset{\handwritten{\scriptscriptstyle\wedge}}{\otimes} U$ : Tordant par $L$ et utilisant le fait que l'on a une compatibilité avec $\otimes$, on trouve un pro-anneau $V$ sur $k'$, et ce dernier hérite de $U$ une application diagonale \ldots
\margin{?}
% left margin, a), alongside source lines 9--16
\margin{\handwritten{a) donner un sens intrinsèque à $C\otimes_k k'$, indépendamment de \deleted{$\unsure{\otimes|\circ}$} la donnée de $F$ (\uline{NB} si $k'$ n'est pas un corps, ce n'est pas une catégorie \uline{abélienne}, seulement additive)}}
% source lines 15--29
Ceci\deleted{x} dit, si on a une donnée $\handwritten{\sigma\in}U\overset{\handwritten{\scriptscriptstyle\wedge}}{\otimes} U$ satisfaisant les conditions (1)\inserted{plus haut,} peut-on définir en termes de $\handwritten{\sigma}$ (ou existe-t-il tout au moins) un élément $\handwritten{\tau\in}V\overset{\handwritten{\scriptscriptstyle\wedge}}{\otimes}_k V$ satisfaisant les relations analogues à (1), de telle façon que $(G,c)$ soit compatible aux isomorphismes de symétrie $\handwritten{\psi}$ dans $C$ défini par $\handwritten{\sigma}$, $\handwritten{\psi'}$ dans $C'$ défini par $\handwritten{\tau}$ ??? La seule chance de \deleted{\unsure{montrer}} construire un $\handwritten{\tau}$ canoniquement semble de prouver que $\handwritten{\tau}$ est ``invariant'' lorsqu'on fait opérer $U$ sur $U\overset{\handwritten{\scriptscriptstyle\wedge}}{\otimes}_k U=(U,\mathrm{int})\overset{\handwritten{\scriptscriptstyle\wedge}}{\otimes}(U,\mathrm{int})$ grâce à son \deleted{\unsure{s}}action dite plus haut sur les deux facteurs, ce qui s'exprime par la relation $U^+ .\handwritten{\tau}=0$.\inserted{\handwritten{[Vérifier dans le cas gradué ce qui en est !]}} - Si ceci marchait (ce qui améliorerait le résultat souhaité en NB au haut de la page), cela montrerait que dans l'étude des foncteurs fibres, il y a lieu de laisser tomber les questions de compatibilité aux symétries, celles-ci se récupérant toujours après coup.\inserted{\handwritten{(Mais cf exemples n\textsuperscript{o}12).}}
% left margin, b), alongside source lines 17--21
\margin{\handwritten{b) Montrer que structure $\otimes$, éventuellement avec données $\varphi,\psi,1,\vee$. se prolonge automatiquement à $C'$, avec \unsure{ses axiomes}}}
% left margin, c), alongside source lines 22--30
\margin{\handwritten{c) généraliser, pour un foncteur comme $G:\unsure{C|\mathcal{C}}\to\mathrm{Mod}(k')$, la description du n\textsuperscript{o}6 en description de $C'$ en termes d'un pro-anneau $U'=\mathrm{End}(G)$, et expliciter comment on a encore les diverses structures sur $C$: $\otimes,\varphi,\psi,1,\vee$}}

% source lines 30--31; word continues on scan 130
Il resterait alors la question d'une description \uline{intrinsèque} d'une catégorie $C$ comme dans le n\textsuperscript{o}1,\inserted{\handwritten{avec $\unsure{\check{X}|\check{\otimes}}$}} \unsure{indépendamment|indépendammment} de la donnée d'un $F$ préli%
\scanpage{130}%
% archive sheet number: 129
% source lines 1--2; completes préliminaire from scan 129
minaire, - sous réserve tout au plus qu'il existe un tel $\unsure{\mathbb{E}|\mathbb{F}|E|F}$ sur une extension convenable de $k_x$ (chose peut-être toujours vrai\handwritten{e}).

% source lines 3--12
12. \uline{Exemples instructifs}. Soit $G$ un monoïde fini. Contrairement au yoga galoisien général, au lieu de prendre $A=k(G)$ (anneau des fonctions sur l'ensemble sous-jacent à $G$) prenons $U=k(G)$, avec l\deleted{a}'application diagonale déduite de $GxG\handwritten{\longrightarrow}G$ (la composition). (On a la bialgèbre ``\unsure{duale|daale}''\inserted{au sens Cartier} de celle qu'on considère d'habitude). Donc la catégorie $C$ définie par $G$ \deleted{est} peut être considérée comme la catégorie des familles $(\unsure{\mathbb{M}|\mathbb{X}}_g)_{g\ G}$ de vectoriels indexés par $G$, \unsure{lé|le}\deleted{application diagonale} produit\inserted{tens.} de deux tels $X$ et $Y$ étant $Z$ défini par $Z_g=\handwritten{\displaystyle\sum}_{g'g''=g}X_{g'}\otimes Y_{g''}$. L'unité est la ``famille de Dirac'' définie par $\mathbb{M}_g=0$ si $g\handwritten{\ne}\mathbb{e}$, $\mathbb{M}_e=k$.\inserted{On a une \unsure{Catégorie|catégorie} avec $\otimes,\handwritten{\phi},1$, et de plus un $\handwritten{\psi}$ si $G$ est commutatif, un $\vee$ si $G$ est un groupe.}

% source lines 13--24, with one displayed formula and its right marginal note
Un $L$ avec une application diagonale $L\handwritten{\longrightarrow}L\otimes_k L$ compatible avec $U\handwritten{\to}U\otimes_k U$ et donnant un isomorphisme par extensions des anneaux de base, peut s'interpréter comme une famille $(\deleted{M}\handwritten{L}_g)_{g\handwritten{\in}G}$, avec des isomorphismes donnés
\[
\deleted{M}\handwritten{L}_{gg'}\handwritten{\simeq}\deleted{M}\handwritten{L}_g\otimes\deleted{M}\handwritten{L}_{g'}\ .
\tag*{(*)}
\]
\margin{Ceci implique que le rang des $\deleted{M}\handwritten{L}_g$ est $0$ ou $1$}
C'est compatible avec $1$ si on a un isomorphisme\inserted{\handwritten{\unsure{(**)|(*+)}}} $\deleted{M}\handwritten{L}_e\handwritten{\simeq}k$, qui rendre les isomorphismes (*) triviaux quand \deleted{\unsure{ceux-ci}} $g$ ou $g'$ est $=1$. C'est associatif si et seulement si les isomorphismes (*) satisfont la \deleted{\unsure{conditio}} conditions d'associativité évidente. De même, $G$ étant commutatif, pour la commutativité. - Le foncteur $F'$ est toujours exact, car $M$ est plat sur $U$.

% source lines 25--30
a) Exemple d'un $\otimes$-foncteur exact, compatible avec $\handwritten{\phi,\psi},1$, \uline{mais non fidèle} (donc pas loc. \unsure{isoM.|isom.} au $\otimes$-foncteur donné). Ce n'est possible dans notre exemple que si $G$ n'est pas un groupe (\deleted{\unsure{sinon $N$|sinon $M$}} car $\deleted{M}\handwritten{L}_g$ est de rang $1$ pour $g$ inversible). Mais prenons $G=(e,p)$, avec $p^{\handwritten{2}}=p$, prenons $\deleted{M}\handwritten{L}_e=k$, $\deleted{M}\handwritten{L}_p=0$. On n'a plus le choix pour les isomorphismes $*$, qui satisfont toutes les compatibilités (avec $1_x$ par construction, $\handwritten{\phi,\psi}$).
% bottom margin, below a handwritten separating rule
\margin{\handwritten{\uline{Question} Un $\otimes$-foncteur compatible avec tout entre $\otimes$-catégories satisfaisant tous les axiomes (y compris rigidité) est-il néc. fidèle et exact ?!}}
